\documentclass[10pt,a4paper]{amsart}
\usepackage[margin=19mm]{geometry}
\usepackage{amsmath,amssymb,mathtools}
\usepackage[scr=rsfs,cal=euler]{mathalfa}
\usepackage{xfrac}
\usepackage[unicode,hidelinks,bookmarks=false]{hyperref}
\newcommand{\ie}{i.e.\ }
\newcommand{\cf}{cf.\ }
\newcommand{\resp}{respectively\ }
\newcommand{\Subsection}{\subsection}
\providecommand{\nobreakdash}{\nobreak\hbox{-}}
\setlength{\emergencystretch}{2em}
\multlinegap0pt
\usepackage{amssymb}
\usepackage[shortlabels]{enumitem}
\newtheorem{lemma}{Lemma}
\newcommand{\Q}{\mathbb{Q}}
\newcommand{\R}{\mathbb{R}}
\newcommand{\Z}{\mathbb{Z}}
\newcommand{\rc}{\operatorname{rc}}
\newcommand{\zero}{\mathbf{0}}
\title{The relaxation complexity of the standard simplex is logarithmic}
\author{Gennadiy Averkov, Simon Keil, Stefan Weltge}
\date{Source: arXiv:2606.11852v2 (CC BY 4.0)}
\begin{document}
\maketitle

\noindent\emph{Source and licence.} The relaxation complexity of the standard simplex is logarithmic. Version arXiv:2606.11852v2 (CC BY 4.0), distributed by arXiv under the Creative Commons Attribution 4.0 International licence, http://creativecommons.org/licenses/by/4.0/, as recorded in the arXiv licence metadata for this exact version and shown on the abstract page https://arxiv.org/abs/2606.11852v2. Full licence notice: \texttt{licenses/arxiv-2606.11852-CC-BY-4.0.txt}.\par\smallskip

\noindent\textit{Editorial scope.} Counted unit: the article's lemma lemFreeJoinRelaxation (source0.tex lines 100--103). The article's complete original proof of it is delivered with it (source0.tex lines 104--124, one \texttt{proof} environment of 21 lines). No mathematics is written, repaired or re-derived: the statement and the proof are the article's own lines, in the article's own notation, and no retained line is substituted or re-flowed.  No further source-local context is needed: every cross-reference of every retained line resolves inside the two delivered spans.  Nothing was omitted from any retained span, so the package carries no omission record.  No retained line contains a citation, so the wrapper carries no bibliography and no bibliography entry.  No retained line contains a figure environment, an image-inclusion command or drawing code, so no image is delivered and the package declares no asset dependency.  Editorial formatting only, in addition: an AMS \texttt{amsart} preamble that restates the article's own declarations and macros used by the retained lines, namely the environments \texttt{lemma} on separate counters, together with the standard packages those lines need.  The retained environments print in this document as Lemma 1 in order of appearance, whereas the article prints the same items with the numbering of its own full document.  The complete native source of the article as distributed in the arXiv e-print for this exact version is bundled unchanged as \texttt{sources/202579/source0.tex}.
\begin{lemma}
    \label{lemFreeJoinRelaxation}
    If $X \subseteq \Z^d$ is finite, then $\rc(X \ast X ) \le \rc(X) + 4$.
\end{lemma}

\begin{proof}
    Let $P$ be a relaxation of $X$ and let $\alpha_1, \dots, \alpha_d \in \R$ be linearly independent over $\Q$.
    Define $g : \R^d \to \R$ by $g(x) = \sum_{i=1}^d \alpha_i x_i$ and set $M = \max_{x \in X} |g(x)|$.
    We claim that the polyhedron $Q$ given by
    \[
        Q = \{(x,y,h) \in \R^d \times \R^d \times \R : |g(x)| \le M (1-h), |g(y)| \le M h, x + y \in P\}
    \]
    is a relaxation of $X \ast X$.
    Since $Q$ has at most 4 more facets than $P$, this implies the statement of the lemma.

    Let $(x,y,h) \in X \ast X$. If $h = 0$, then $y = \zero$ and $x \in X$.
    By the choice of $M$, this implies $(x,y,h) \in Q$. 
    If $h = 1$, then $x = \zero$ and $y \in X$, and we again obtain $(x,y,h) \in Q$.
    Thus, $X \ast X \subseteq Q$.

    To see that $Q \cap \Z^{2d+1} \subseteq X \ast X$ holds, let $(x,y,h) \in Q \cap \Z^{2d+1}$.
    Since $|g(x)|, |g(y)|$ and $M$ are non-negative, the inequalities imply $0 \le 1 - h$ and $0 \le h$, which together with the integrality gives $h \in \{0,1\}$. If $h = 0$, then $|g(y)| \le M h$ implies $g(y) = 0$.
    Since the $\alpha_i$ are linearly independent over $\Q$ and $y$ is integer, this implies $y = \zero$.
    Thus, we have $x = x + y \in P$ and hence $(x,y,h) = (x, \zero, 0) \in X \ast X$.
    If $h = 1$, we analogously obtain $x = \zero$ and $y \in X$, which implies $(x,y,h) = (\zero, y, 1) \in X \ast X$.
\end{proof}

\end{document}
